\section{Introduction}
\label{sec:introduction}

\subsection{Context and Motivation: The Need for Architectural Code Review}
\label{sec:intro:context}

Automated code review and software quality evaluation have become indispensable pillars of modern software engineering \cite{tufano2022towards}. As distributed software systems scale in complexity and release cycles compress under continuous integration and continuous delivery (CI/CD) regimes, automated tools must detect regressions, evaluate non-functional requirements, and recommend actionable remediations before defective changes reach production fabrics.

Distributed publish--subscribe middleware frameworks---including the Robot Operating System (ROS~2), the Object Management Group's Data Distribution Service (DDS), and MQTT---form the communication infrastructure of contemporary microservices, IoT smart cities, and safety-critical cyber-physical systems \cite{lee2025probabilistic, lee2025dependency}. These frameworks decouple message producers and consumers across space, time, and synchronization via topic-based publish--subscribe channels and message brokers. However, this architectural decoupling introduces an insidious challenge for automated quality assurance: it obscures the end-to-end dependency structure of the system. Failures, buffer overflows, and latency spikes propagate along indirect topological pathways that are completely invisible when analyzing source code files in isolation.

\subsection{Two Open Gaps: The Architecture-Code Gap and Open-Loop Review}
\label{sec:intro:gaps}

Quality evaluation in continuous integration has historically relied on Static Code Analysis (SCA) platforms (e.g., SonarQube, SpotBugs, ESLint). This creates a fundamental \textbf{Architecture-Code Gap}:
\begin{enumerate}
    \item \textbf{The Architecture-Code Gap:} A distributed system can achieve pristine code quality within every individual component---clean class hierarchies, zero file-level code smells, and low cyclomatic complexity---yet remain brittle at the architectural topology level. Vulnerabilities such as single points of failure (SPOFs), severe topic fan-out bottlenecks, co-located deployment risks, and incompatible Quality of Service (QoS) transport contracts cannot be detected by examining source files in isolation. Shifting structural dependability checks left into the development lifecycle demands a transition from traditional Static Code Analysis to \emph{Static System Analysis (SSA)}.
    \item \textbf{The Unnamed-Pathology and Open-Loop Review Gap:} In object-oriented programming, mature catalogs of bad smells and anti-patterns (e.g., God Class, Feature Envy, Shotgun Surgery) provide developers with a shared vocabulary, testable detection rules, and proven refactoring patterns \cite{fowler1999refactoring, brown1998antipatterns, suryanarayana2014refactoring}. Microservices research has similarly begun cataloging smells (e.g., distributed monoliths, cyclic dependencies, shared databases) \cite{taibi2020microservices}. In contrast, \emph{publish--subscribe architectures lack an equivalent formal anti-pattern catalog}. Consequently, architectural flaws are discovered reactively through production cascade postmortems rather than proactively during code review. Moreover, existing topology-aware diagnostic frameworks typically operate \emph{open-loop}: they compute abstract numerical centrality scores or black-box failure-impact predictions without identifying the specific architectural pathology at fault or synthesizing verified, actionable refactoring guidance \cite{harman2012sbse, aleti2013software}.
\end{enumerate}

\subsection{The Diagnostic Pathway: Transparent, Explainable Software Quality Evaluation}
\label{sec:intro:solution}

To address these gaps, this paper brings the \textbf{Diagnostic Pathway} to the forefront of automated code review and software quality evaluation in distributed systems. We present \textbf{SaG-Prescribe}, a comprehensive framework that integrates deterministic diagnostic quality attribution, anti-pattern detection, and closed-loop prescriptive refactoring within continuous integration workflows.

SaG-Prescribe decouples quality assurance into two principled, cooperative pathways:
\begin{itemize}
    \item \textbf{The Diagnostic Pathway (Deterministic Software Quality Evaluation \& Attribution):} The Diagnostic Pathway operates statically on Architecture-as-Code descriptors. It constructs a multi-layer heterogeneous dependency graph ($G_{\text{analysis}}$) encompassing applications, libraries, topics, brokers, and physical hosting nodes. It bridges code-level and architecture-level quality by synthesizing modular SCA metrics into a \textbf{Code Quality Penalty (CQP)}. It then applies a hierarchical \textbf{Reliability--Maintainability (RM)} attribution model grounded in the \textbf{ISO/IEC 25010} (Product Quality) and \textbf{ISO/IEC 25019} (Quality-in-Use) standards \cite{iso25010_2023, iso25019_2023}. Reliability decomposes into \textbf{Fault Tolerance} ($FT$, error propagation reach) and \textbf{Availability} ($A$, structural single-point-of-failure exposure), while \textbf{Maintainability} ($M$) quantifies coupling complexity and change ripple risk. The diagnostic pathway maps structural technical debt directly to stakeholder harm categories ($\mathbf{h}_{\text{QiU}}$) and routes explainable findings to specific engineering roles (Reliability Engineers, DevOps/SREs, and Software Architects). Over these metric vectors, it executes a catalog of nineteen formal publish--subscribe anti-patterns using adaptive box-plot thresholds.
    \item \textbf{The Counterfactual Prescriptive Pathway (Closed-Loop Refactoring Verification):} Rather than offering open-loop suggestions, the prescriptive engine compiles diagnostic findings into three typed graph-mutation operators: logical topic splitting, physical anti-affinity reallocation, and transport QoS contract hardening. It subjects every candidate refactoring to \textbf{counterfactual verification} on an in-memory sandbox graph (\texttt{MemoryRepository}), simulating failure cascades across multiple seeds and propagation thresholds. An edit is admitted if and only if its measured impact reduction exceeds the simulator's seed noise at every threshold ($\overline{\Delta I} > \kappa \sigma_{\text{seed}}$), and the composite policy is gated on a net positive System Risk Index delta ($\Delta\text{SRI} > 0$).
\end{itemize}

\subsection{Contributions}
\label{sec:intro:contributions}

This paper makes the following contributions to automated software quality assurance and code review:
\begin{enumerate}
    \item \textbf{A Formal Diagnostic Pathway for Pub-Sub Software Quality Evaluation (Section~\ref{sec:sag_model}):} We formalize a multi-level diagnostic model that ingests file-level SCA metrics via the Code Quality Penalty (CQP), derives multi-layer dependency projections ($G_{\text{analysis}}$), and computes a hierarchical Reliability--Maintainability (RM) attribution model grounded in ISO/IEC 25010 and ISO/IEC 25019 Quality-in-Use standards.
    \item \textbf{A Comprehensive Catalog of 19 Publish--Subscribe Anti-Patterns and Bad Smells (Section~\ref{sec:antipattern_catalog}):} We formalize nineteen architectural anti-patterns and smells as topological signatures over a 16-element structural metric vector with adaptive box-plot thresholds, providing explainable review feedback for automated review bots.
    \item \textbf{A Prescriptive Refactoring Pipeline with Counterfactual Verification (Sections~\ref{sec:optimization_objective}--\ref{sec:prescriptive_pipeline}):} We introduce a rule-based refactoring compiler translating diagnostic findings into three graph-mutation operators (topic splitting, anti-affinity reallocation, QoS hardening), verified against a discrete-event cascade simulator using a two-level margin-gating criterion. We disclose the exact automation footprint (5 of 19 patterns automated, 14 advisory).
    \item \textbf{An Automated Code Review Bot and CI/CD Quality Gate (Section~\ref{sec:cicd_gating}):} We operationalize the Diagnostic Pathway as a sub-second pre-merge quality gate featuring standardized exit codes (0/1/2) and formalize delta-aware merge-base regression semantics.
    \item \textbf{Empirical Evaluation and Methodological Findings (Sections~\ref{sec:experimental_design}--\ref{sec:results}):} Evaluated across eight benchmark scenarios (12 to 520 components), we report: (a)~diagnostic detection achieves mean Spearman $\rho = 0.485$ against simulated cascade impact, functioning as a sensitive screen ($\text{recall} = 0.942$, $\text{precision} = 0.253$, implicating $94.8\%$ of components) and exhibiting a Simpson's paradox pooling effect ($\rho = 0.028$ pooled vs. $0.503$ on applications); (b)~counterfactual verification admits 1128 of 1589 candidates ($71.0\%$), with anti-affinity reallocation demonstrating $77.7\%$ survival, and every scenario achieving statistically significant risk reduction (Wilcoxon $W=0, p=0.0156, n=7$); (c)~diagnostic review gating executes in 0.01--20.98~s, confirming real-time CI/CD viability.
\end{enumerate}

\subsection{Relationship to Prior Work}
\label{sec:intro:prior_work}

This submission builds upon the foundational graph schema and simulation infrastructure introduced in our companion work \cite{sag_companion}, while contributing novel, unshared capabilities:
\begin{itemize}
    \item Companion manuscript \cite{sag_companion} introduces the heterogeneous graph schema, simulation kernels, and a learned Graph Neural Network (GNN) for failure-impact prediction.
    \item This paper contributes the \textbf{Diagnostic Pathway} for automated code review, the Code Quality Penalty (CQP), the 19-pattern pub-sub anti-pattern catalog, the prescriptive refactoring compiler, the counterfactual verification engine, the CI/CD quality gate, and the comprehensive empirical evaluation of detection and prescription.
    \item Detection figures previously reported from preliminary conference workshops are re-measured here with strict artifact reproducibility, reporting negative results and harness corrections with complete scientific candor.
\end{itemize}

\subsection{Organization}
\label{sec:intro:org}

Section~\ref{sec:related_work} surveys related work. Section~\ref{sec:sag_model} formalizes the Diagnostic Pathway and software quality model. Section~\ref{sec:antipattern_catalog} presents the anti-pattern catalog. Section~\ref{sec:optimization_objective} defines the closed-loop optimization objective. Section~\ref{sec:prescriptive_pipeline} details the prescriptive pipeline. Section~\ref{sec:cicd_gating} describes CI/CD integration. Sections~\ref{sec:experimental_design} and \ref{sec:results} present the experimental design and results. Section~\ref{sec:discussion} discusses implications and threats to validity, and Section~\ref{sec:conclusion} concludes.
